\section{Objetivos}
\subsection{Objetivo General}
\begin{itemize}

	\item Comparar el desempeño computacional de una técnica de visualización de campos de flujo (Streamlines, Streaklines y Pathlines), que involucra el método de Runge-Kutta implementado en lenguaje C/C++, mediante métodos tradicionales de cómputo y programación heterogénea con SYCL, evaluando métricas como el tiempo de ejecución, uso de memoria y escalabilidad.

\end{itemize}

\subsection{Objetivos Específicos}

\begin{itemize}
	\item Formular una versión optimizada del algoritmo RK4 en SYCL, con el propósito de evaluar su desempeño en arquitecturas heterogéneas.
	\item Comparar el comportamiento de las implementaciones en C/C++ y en SYCL bajo diferentes configuraciones del problema (tamaño del dataset, conjunto de semillas y complejidad del campo vectorial), con el fin de identificar diferencias en su eficiencia y desempeño.
	\item Analizar comparativamente las ventajas, limitaciones y condiciones de uso de SYCL frente a implementaciones tradicionales en CPU y en GPU con CUDA, considerando métricas de tiempo de ejecución, uso de memoria y escalabilidad.
\end{itemize}
